\documentclass[12pt,a4paper]{report}

% --- PACKAGES ---
\usepackage[french]{babel}
\usepackage[T1]{fontenc}
\usepackage[utf8]{inputenc}
\usepackage{listings}
\usepackage{xcolor}
\usepackage[most]{tcolorbox}
\usepackage{geometry}

\geometry{margin=2.5cm}

% --- COULEURS ---
\definecolor{sep-gray}{rgb}{0.4, 0.4, 0.4}
\definecolor{code-bg}{rgb}{0.98, 0.98, 0.98}

% --- STYLE DES BOITES (Design corrigé pour Yahya) ---
\newtcolorbox{sepbox}[1]{
    enhanced,
    attach boxed title to top left={yshift=-2mm, xshift=2mm},
    colback=white,
    colframe=sep-gray,
    fonttitle=\bfseries,
    colbacktitle=sep-gray,
    title=#1,
    boxed title style={size=small, sharp corners},
    drop shadow,
    bottom=2mm
}

% --- CONFIGURATION C ---
\lstset{
    language=C,
    backgroundcolor=\color{code-bg},
    basicstyle=\ttfamily\footnotesize,
    keywordstyle=\color{blue},
    commentstyle=\color{gray},
    stringstyle=\color{purple},
    breaklines=true,
    frame=none,
    numbers=left,
    numberstyle=\tiny\color{gray}
}

\begin{document}

\chapter{Module de Séparation Visuelle (Separators)}

\section{Introduction}
L'organisation visuelle est primordiale pour guider l'œil de l'utilisateur. Le module \texttt{separator} permet d'insérer des lignes de démarcation horizontales ou verticales entre les différents blocs de notre application de transport.

\section{Définition de la Configuration}

\begin{sepbox}{Structure : separator\_config}
Cette structure simple permet de définir l'orientation de la ligne de séparation. Elle hérite également du style commun via \texttt{widget\_style\_config}.
\end{sepbox}

\begin{lstlisting}
typedef struct {
    GtkOrientation orientation;
    widget_style_config style;
} separator_config;
\end{lstlisting}

\section{Macro d'Initialisation par Défaut}

\begin{sepbox}{Macro : SEPARATOR\_CONFIG\_DEFAULT}
Par défaut, le séparateur est configuré en mode horizontal ($GTK\_ORIENTATION\_HORIZONTAL$), ce qui est le cas d'utilisation le plus fréquent pour séparer des sections de texte ou des boutons.
\end{sepbox}

\begin{lstlisting}
#define SEPARATOR_CONFIG_DEFAULT ((separator_config){ \
    .orientation = GTK_ORIENTATION_HORIZONTAL, \
    .style = WIDGET_STYLE_CONFIG_DEFAULT, \
})
\end{lstlisting}

\section{Implémentation de la Création}

\begin{sepbox}{Fonction : create\_separator}
La fonction vérifie si une configuration est fournie. Si c'est le cas, elle crée le widget avec l'orientation demandée et applique le style personnalisé.
\end{sepbox}

\begin{lstlisting}
GtkWidget *create_separator(const separator_config *config) {
    GtkWidget *separator = gtk_separator_new(config ? 
        config->orientation : GTK_ORIENTATION_HORIZONTAL);
    
    if (config != NULL) {
        apply_widget_style(separator, &config->style);
    }
    return separator;
}
\end{lstlisting}

\end{document}